\subsection{Reproducing the closed study}
\label{sec:repro}

The study we build on reported, on its untouched final holdout (7 scenarios
$\times$ seeds 1001--1005, three training roots), a mean return of
\HistBandit{} for the bandit and \HistPPO{} for PPO, and a small PPO advantage
in the deceptive-local-optimum scenario. Its checkpoints were not archived,
so we retrained all six runs with the governed trainer at the same commit
lineage on a CPU-only machine (the originals used CUDA) and repeated its
selection and holdout procedure exactly.

Three properties reproduce \emph{exactly}: the \BaselineReproCells{} baseline (policy,
scenario) holdout means and standard deviations (to $\BaselineReproMaxDiff$), the
training episode-seed ledger of every root (byte-identical), and all integrity
counters. The bandit reproduces statistically (\ReproBandit{} vs.\
\HistBandit{}; Table~\ref{tab:repro}); PPO does not (\ReproPPO{} vs.\
\HistPPO{}). With identical training traffic, the executions differ in device
(CUDA vs.\ CPU), in library versions (PyTorch, NumPy, Python patch level) and
hence in floating-point arithmetic and in PPO's action-sampling random stream,
which SB3 draws from the device's generator; the bandit's exploration and
replay sampling use a NumPy generator and are device-independent. These
differences moved PPO's per-root holdout return by up to \ReproMaxShiftPPO{}
points and the bandit's by at most \ReproMaxShiftBandit{}. The bandit is not
immune to such changes: on root 42, changing only its learner seed moves its
test return by up to \NoiseBanditSpread{} (\S\ref{sec:rq3}). The comparison is
therefore not a pure numerical-perturbation experiment for PPO; it shows that
PPO's outcome depends strongly on incidental properties of an execution, which
the spread across our five CPU training roots confirms (root standard
deviation \TestPPORootSd{} for PPO vs.\ \TestBanditRootSd{} for the
bandit). On
\ReproPPONeverPositive{} of 3 roots the reproduced PPO never reached a positive
validation return at any checkpoint. The bandit's advantage reappeared on all
three roots and grew to \ReproGap{} (per root \ReproGapPerRoot{}, roots 42 /
271828 / 314159; with three roots the $t$-interval [\ReproGapTLo,
\ReproGapTHi] is uninformative about magnitude); the historical deceptive-scenario PPO advantage
(\HistDeceptive{} for the bandit) did not reproduce (bandit $-$ PPO
$=\ReproDeceptive{}$, positive on all three roots and in every paired
episode).

\begin{table}[t]
\centering\small
\caption{Closed-study reproduction under its own protocol (selection on seeds
101--105, holdout 1001--1005, 35 episodes). Selected checkpoint and mean
holdout return, historical (CUDA) vs.\ reproduced (CPU).}
\label{tab:repro}
\begin{tabular}{@{}lrcrr@{}}
\toprule
Learner & Root & Selected (hist./repr.) & Hist. & Repr. \\
\midrule
\tabinput{tables/reproduction_per_root}
\bottomrule
\end{tabular}
\end{table}

We draw two methodological conclusions before any new experiment. First,
PPO's outcome in this environment has a large run-to-run variance that is not
captured by three roots, so the remainder of the paper uses five roots for
the primary comparison and reports root-level uncertainty. Second, the
bandit's result is robust to the perturbation; this is the first indication
that its advantage is not an accident of one execution.
